\documentclass{mudasir}


\date{}
\author{Mudasir}

\begin{document}

\title{Examples}
\maketitle

This file demonstrates the features provided by the mudasir class. Most examples are AI generated as a simple demo of the visual style of the various commands.

\section*{Math Macro Examples}
\begin{align*}
    \sni{1} \frac{1}{n^2} = \frac{\pi^2}{6} \\
    \df{f(x)}{\int_0^x f(t) dt} = 1 \\
    \cancelto{0}{\int_0^x f(t) dt}
\end{align*}

Derivatives:
\[
    \deriv{y}{x},\quad \deriv{e^{x^2}}{x}=2x e^{x^2}
\]

Partial derivatives:
\[
    \pderiv{u}{x},\quad \pderiv{^2 u}{x^2}
\]

Integrals:
\[
    \integral{0}{1}{x^2}{x} = \left.\frac{x^3}{3}\right|_0^1 = \frac{1}{3}
\]

Differentials (explicit \verb|d|):
\[
    \int f(t)\,\diff{t}
\]

Vector calculus / del operator:
\[
    \del f,\qquad \del\cdot \vec{F}
\]

Laplace transform:
\[
    \laplace{f(t)} = F(s) = \int_0^{\infty} f(t)e^{-st}\,\diff{t}
\]

Inline code example: \code{disp('Hello, world!')}

\section*{Theorem-like Environments}
\begin{theorem}
Let $x\in\RR$. Then $x^2\ge0$.
\end{theorem}

\begin{theorem*}
The square of any real number is nonnegative (unnumbered).
\end{theorem*}

\begin{lemma}
If $a>0$ and $b>0$ then $ab>0$.
\end{lemma}

\begin{lemma*}
An unnumbered lemma example.
\end{lemma*}

\begin{proposition}
Proposition: addition of reals is commutative.
\end{proposition}

\begin{corollary}
Corollary: immediate consequence of the proposition.
\end{corollary}

\section*{Definitions, Claims, and Examples}
\begin{definition}
Definition of continuity at a point: A function $f$ is continuous at $x=a$ if $\lim_{x\to a} f(x) = f(a)$.
\end{definition}

\begin{definition*}
An unnumbered definition.
\end{definition*}

\begin{claim}
This is a sample claim with a short justification.
\end{claim}

\begin{claim*}
Unnumbered claim demonstration.
\end{claim*}

\begin{example}
Example: compute $1+1=2$.
\end{example}

\section*{Problems and Exercises}
\begin{problem}
Compute the integral $\int_0^1 x^2\,dx$.
\end{problem}

\begin{niceprob}
A highlighted problem with a star symbol.
\end{niceprob}

\begin{exercise}
Exercise: Show that the derivative of $e^{x}$ is $e^{x}$.
\end{exercise}

\begin{exercise*}
Unnumbered exercise: verify a trivial identity.
\end{exercise*}

\section*{Notes, Remarks, Questions, Answers}
\begin{note}
This is a side note (unnumbered by style).
\end{note}

\begin{remark}
Remark: pay attention to edge cases.
\end{remark}

\begin{remark*}
An unnumbered remark.
\end{remark*}

\begin{question}
What is the limit of $\sin x / x$ as $x\to0$?
\end{question}

\begin{question*}
An unnumbered question.
\end{question*}

\begin{answer}
Answer: $1$.
\end{answer}

\begin{answer*}
An unnumbered answer entry.
\end{answer*}

\begin{reminder}
Reminder: check assumptions in proofs.
\end{reminder}

\begin{reminder*}
Unnumbered reminder.
\end{reminder*}

\begin{warning}
Warning: don't divide by zero.
\end{warning}

\begin{warning*}
Unnumbered warning.
\end{warning*}

\section*{Solution, Walkthrough, and Custom List}
\begin{problem}
Prove that $2+2=4$.
\begin{soln}
Trivial by arithmetic.
\end{soln}
\end{problem}

\begin{walkthrough}
This is a brief walkthrough for a sample solution.
\end{walkthrough}

\begin{walk}
\ii First step.
\ii Second step.
\ii Third step.
\end{walk}

\section*{Code (minted)}
Block code example (from `test.m`):
\inputminted[frame=single, linenos]{matlab}{test.m}

\end{document}